\appendix
\customchapter{ANEXOS}

\section*{Anexo A. Configuración de \emph{software}}

Python 3.10, PyTorch, MinkowskiEngine, \texttt{spconv}, Pointcept, código oficial
de HyperNCD y pesos públicos de Sonata. Seguimiento de experimentos con
Weights~\&~Biases.

\section*{Anexo B. Nomenclatura}

\begin{table}[H]
\centering
\setstretch{1}\small
\begin{tabular}{@{}p{2.2cm}p{7.3cm}p{3.2cm}@{}}
\toprule
\textbf{Símbolo} & \textbf{Significado} & \textbf{Valor} \\
\midrule
$P$ & Nube de puntos del levantamiento & Catedral de Lima, TLS \\
$N$ & Número de puntos & $\sim$2$\times$10$^{7}$ tras voxelizar \\
$\mathcal{C}_k$ & Conjunto de clases conocidas (etiquetadas) & Por sección \\
$\mathcal{C}_n$ & Conjunto de clases novedosas (sin etiqueta) & Hasta 6 por sección \\
$f_{\mathrm{sem}}$ & Descriptor semántico (MinkowskiUNet-34C) & 96 dimensiones \\
$f_{\mathrm{geo}}$ & Descriptor geométrico calculado a mano & 6 dimensiones \\
$f_{\mathrm{ssl}}$ & Descriptor auto-supervisado (Sonata, congelado) & 1088 dimensiones \\
$F_i$ & Descriptor compuesto del punto $i$ & Concatenación \\
$K$ & Número de prototipos & Igual al de clases \\
$S_p$ & Matriz de asignación suave & $N\!\cdot\!P\times K$ \\
$S_g$ & Afinidad geométrica entre prototipos & Núcleo gaussiano \\
$S_s$ & Afinidad semántica entre prototipos & Coseno \\
$S_b$ & Afinidad combinada & Ec.~\eqref{eq:afinidad} \\
$\alpha$ & Balance geométrico-semántico en $S_b$ & Inicia en 0{,}5 \\
$M$ & Vecinos por hiperarista & 8 \\
$\tau$ & Temperatura del núcleo gaussiano & 0{,}1 \\
$\epsilon$ & Suavizado de etiquetas & 0{,}15 \\
$G$ & Brecha cerrada & Ec.~\eqref{eq:G} \\
\bottomrule
\end{tabular}
\end{table}

\section*{Anexo C. Glosario de términos}

\begin{center}
\setstretch{1}\small
\begin{longtable}{@{}p{4.0cm}p{8.7cm}@{}}
\toprule
\textbf{Término} & \textbf{Definición} \\
\midrule
\endfirsthead
\toprule
\textbf{Término} & \textbf{Definición} \\
\midrule
\endhead
\bottomrule
\endlastfoot
Afinidad & Semejanza entre dos prototipos; combina una componente geométrica y una semántica. \\
Aprendizaje auto-supervisado & Entrenamiento cuya señal de supervisión procede de los propios datos y no de etiquetas humanas. \\
Atajo geométrico & Colapso de la representación sobre pistas geométricas simples en lugar de aprender semántica. \\
Clase conocida & Clase con puntos etiquetados durante el entrenamiento. \\
Clase novedosa & Clase presente en la nube pero sin ninguna etiqueta. \\
Codificador congelado & Red usada solo en inferencia, cuyos pesos no se actualizan. \\
Descriptor & Vector numérico que resume el contenido de un punto o de una región. \\
Emparejamiento húngaro & Asignación óptima entre agrupaciones descubiertas y clases reales. \\
Hiperarista & Conjunto formado por un prototipo y sus $M$ vecinos más afines. \\
Hipergrafo & Estructura en la que una arista puede conectar más de dos nodos a la vez. \\
\emph{Linear probing} & Evaluación de un descriptor congelado mediante un clasificador lineal. \\
Nube de puntos & Conjunto de posiciones tridimensionales medidas con sus atributos. \\
NCD & Descubrimiento de clases novedosas. \\
Prototipo & Vector representante de un grupo de puntos o regiones semejantes. \\
Pseudo-etiqueta & Etiqueta asignada por el propio modelo, sin supervisión humana. \\
Región o superpunto & Agrupación de puntos vecinos tratada como una sola unidad. \\
Rejilla dispersa & Indexación por coordenada entera de vóxel que almacena solo los ocupados. \\
Sección & Partición espacial disjunta del levantamiento. \\
Segmentación semántica & Asignación de una etiqueta de clase a cada punto. \\
Serialización & Ordenamiento de los puntos mediante una curva de llenado de espacio. \\
Sinkhorn-Knopp & Normalización iterativa que equilibra el reparto de pseudo-etiquetas. \\
TLS & Escaneo láser terrestre desde posiciones fijas. \\
\emph{Up-cast} & Propagación inversa del submuestreo hasta la resolución de entrada. \\
Voxelización & Sustitución por un representante único de los puntos de un mismo cubo. \\
\end{longtable}
\end{center}
